\par\medskip\noindent\textbf{Cómo leer la salida.}
{\small\setlength{\tabcolsep}{4pt}\renewcommand{\arraystretch}{1.15}
\begin{xltabular}{\linewidth}{@{}>{\raggedright\arraybackslash}p{0.24\linewidth}>{\raggedleft\arraybackslash}p{0.15\linewidth}>{\raggedright\arraybackslash}X@{}}
\toprule Indicador & Valor & Interpretación \\ \midrule \endfirsthead
\toprule Indicador & Valor & Interpretación \\ \midrule \endhead
\bottomrule \endlastfoot
CP1 & 73,46\,\% & El primer componente explica casi tres cuartos de la varianza de las 30 carteras. \\
CP2, CP3 & 6,29\,\%; 4,01\,\% & Los dos siguientes explican mucho menos; con tres se llega al 83,76\,\%. \\
CP4, CP5 & 2,87\,\%; 2,03\,\% & Aportes marginales: la estructura de covarianzas tiene pocos factores relevantes. \\
Correlación CP1-Mkt & 0,962 & El primer componente es el mercado. \\
Correlación CP2-SMB & $-0{,}801$ & El segundo es (con signo contrario) el factor tamaño. \\
Correlación CP3-Mom / CP3-HML & 0,855 / $-0{,}562$ & El tercero combina momentum y valor (con signo opuesto). \\
$R^2$ de cada componente & 0,989; 0,836; 0,916 & Los cuatro factores explican casi por completo los tres componentes: los factores de Fama-French y Carhart capturan las fuentes estadísticas de covarianza. \\
$R^2$ medio de las 30 carteras & 0,817 / 0,794 & Tres componentes explican la varianza de las carteras un poco mejor que tres factores, como debe ser por construcción. \\
\end{xltabular}}
